% !TeX root = ../main.tex
\documentclass[../main.tex]{subfiles}

\begin{document}

\chapter{Mecánica}

\section{Primera ley de Newton}
Todo cuerpo continúa en su estado de reposo o de movimiento uniforme en línea recta, a menos que una fuerza externa neta actúe sobre él.

La primera ley de Newton establece, en primer lugar, el movimiento como la piedra fundamental de la física, de acuerdo con las ideas de Aristóteles. 

Segundo; introduce el movimiento con velocidad constante como un \emph{movimiento patrón}, así como el metro, el kilogramo, etc., que nos permite detectar la presencia de fuerzas. Esto es: si un cuerpo se mueve en línea recta con velocidad constante, entonces no hay fuerza neta; de lo contrario, sí hay fuerza neta. Así, dado el movimiento de un cuerpo, lo comparamos con el movimiento patrón para detectar fuerzas.

Tercero: concepto de \textbf{fuerza}.

\textbf{Fuerza:} magnitud física que modifica el estado de movimiento de un cuerpo (lo acelera, frena, cambia el sentido, etc.) o bien lo deforma.

Para entender lo que es el concepto de fuerza debemos introducir el concepto de \textbf{interacción}.

\textbf{Interacción:} es la acción de una parte del Universo sobre otra parte del Universo.

\section{Acción en mecánica}
La palabra \textbf{acción} indica que una persona, animal o cosa (material o inmaterial) está haciendo algo sobre otra persona, animal o cosa; es decir, está actuando (de manera voluntaria o involuntaria, de pensamiento, palabra u obra) sobre algo, lo que normalmente implica movimiento o cambio de estado o situación y afecta o influye en una persona, animal o cosa.

En mecánica nos interesa la acción que genera movimiento: una parte 1 del Universo (\textbf{agente externo}) al actuar sobre otra parte 2 del Universo (\textbf{sistema}) produce en 2 un cambio en su estado de movimiento. Por ejemplo, una estrella interactuando sobre un planeta, o una bola de billar en movimiento interactuando sobre otra bola de billar en reposo.



\newpage

\end{document}
